\section{Tối ưu hóa truy xuất dữ liệu}

Dựa trên pipeline xử lý truy vấn của hệ quản trị CSDL, việc viết câu lệnh SQL là chưa đủ. Hệ thống cần tối ưu hóa cây truy vấn (Query Tree) thông qua các luật biến đổi tương đương để giảm thiểu chi phí bộ nhớ và I/O trước khi lập kế hoạch thực thi (Execution Plan).

\subsection{Bài toán truy vấn}
\textbf{Yêu cầu:} Lấy ra Mã giao dịch (\texttt{TransactionID}) và Số điểm (\texttt{Amount}) của các giao dịch cộng điểm lớn hơn 500, thuộc về các Chiến dịch (\texttt{CAMPAIGN}) đang ở trạng thái hoạt động (\texttt{'Active'}).

\textbf{Câu lệnh SQL đại số quan hệ:}
$$\pi_{TransactionID, Amount} ( \sigma_{Amount > 500 \wedge Status = \text{'Active'}} (CREDIT\_TRANSACTION \bowtie CAMPAIGN) )$$

\subsection{Cây truy vấn trước và sau khi tối ưu}

Theo luật biến đổi tương đương, phép chọn (Selection - $\sigma$) nếu được đẩy xuống trước phép kết nối (Join - $\bowtie$) sẽ giúp giảm thiểu đáng kể kích thước dữ liệu trung gian.

\begin{figure}[H]
    \centering
    \resizebox{\textwidth}{!}{ % Ép toàn bộ sơ đồ tự động co lại vừa khít 100% chiều ngang trang giấy
    \begin{tikzpicture}[
        box/.style={rectangle, draw=black, rounded corners, inner sep=12pt, text centered, font=\normalsize},
        leaf/.style={rectangle, draw=black, inner sep=10pt, text centered, font=\normalsize\bfseries}
    ]
        % ================= CÂY CHƯA TỐI ƯU =================
        \node[font=\Large\bfseries] at (0, 8.5) {Chưa tối ưu};
        
        \node[box] (proj1) at (0, 7) {$\pi_{TransactionID, Amount}$};
        \node[box] (sel1) at (0, 5) {$\sigma_{Amount > 500 \wedge Status = \text{'Active'}}$};
        \node[box] (join1) at (0, 3) {$\bowtie_{CampaignID}$};
        
        \node[leaf] (trans1) at (-4.5, 0) {CREDIT\_TRANSACTION};
        \node[leaf] (camp1) at (4.5, 0) {CAMPAIGN};
        
        \draw[->, thick] (sel1) -- (proj1);
        \draw[->, thick] (join1) -- (sel1);
        \draw[->, thick] (trans1) -- (join1);
        \draw[->, thick] (camp1) -- (join1);

        % ================= CÂY ĐÃ TỐI ƯU =================
        \node[font=\Large\bfseries] at (18, 8.5) {Đã đẩy phép chọn};
        
        \node[box] (proj2) at (18, 7) {$\pi_{TransactionID, Amount}$};
        \node[box] (join2) at (18, 5) {$\bowtie_{CampaignID}$};
        
        \node[box] (sel2_trans) at (13, 2.5) {$\sigma_{Amount > 500}$};
        \node[box] (sel2_camp) at (23, 2.5) {$\sigma_{Status = \text{'Active'}}$};
        
        \node[leaf] (trans2) at (13, 0) {CREDIT\_TRANSACTION};
        \node[leaf] (camp2) at (23, 0) {CAMPAIGN};
        
        \draw[->, thick] (join2) -- (proj2);
        \draw[->, thick] (sel2_trans) -- (join2);
        \draw[->, thick] (sel2_camp) -- (join2);
        \draw[->, thick] (trans2) -- (sel2_trans);
        \draw[->, thick] (camp2) -- (sel2_camp);
        
        % Đường phân cách dọc ở giữa
        \draw[dashed, thick, gray] (9, -1) -- (9, 9.5);
    \end{tikzpicture}
    }
    \caption{Cây truy vấn trước và sau khi tối ưu hóa bằng luật phân rã}
\end{figure}

\subsection{Đánh giá thuật toán kết nối (Join Algorithm)}
Bằng việc đẩy phép chọn xuống dưới, kích thước của 2 bảng trước khi thực hiện phép Join đã bị thu hẹp đáng kể. 

Tuy nhiên, thay vì khẳng định hệ thống luôn chọn một thuật toán cố định, ta cần hiểu rằng Trình tối ưu hóa truy vấn (Query Optimizer) của SQL Server sẽ đánh giá dựa trên thống kê thực tế. Nếu bảng \texttt{CAMPAIGN} sau khi lọc trở nên rất nhỏ và có thể nạp toàn bộ vào bộ nhớ đệm (Buffer pool), thuật toán \textbf{Nested-Loop Join} có thể được hệ thống ưu tiên. Ngược lại, nếu kích thước dữ liệu trung gian của cả hai bảng vẫn lớn và không có thứ tự, hệ thống sẽ linh hoạt chuyển sang sử dụng \textbf{Hash Join} để đạt chi phí $O(b_R + b_S)$.